\documentclass[dvisvgm,tikz,border=3pt]{standalone}
\usepackage{amsmath,amssymb}
\usepackage{pgfplots}
\usepackage{CJKutf8}
\pgfplotsset{compat=1.18}
\usetikzlibrary{arrows.meta,calc}

\definecolor{themebg}{HTML}{30362d}
\definecolor{themefg}{HTML}{edf4e8}
\definecolor{themegray}{HTML}{9ea897}
\definecolor{themecurve}{HTML}{7eb6ff}
\definecolor{themealert}{HTML}{ff7b7b}
\definecolor{themeamber}{HTML}{f5b942}

\begin{document}
\begin{CJK*}{UTF8}{gbsn}
\pagecolor{themebg}
\begin{tikzpicture}[color=themefg, text=themefg]
\begin{scope}[x=8.2cm,y=8.2cm]
  % 1. 坐标轴与辅助虚线（底层）
  \draw[themefg, line width=0.9pt, -{Stealth[scale=1.3, length=3.0mm]}]
    (-0.035,0) -- (1.08,0) node[below left=1pt] {$x$};
  \draw[themefg, line width=0.9pt, -{Stealth[scale=1.3, length=3.0mm]}]
    (0,-0.035) -- (0,1.08) node[below left=1pt] {$y$};

  \draw[themegray, densely dotted, line width=0.8pt] (0.5,0) -- (0.5,0.5);
  \draw[themegray, densely dotted, line width=0.8pt] (0,0.5) -- (0.5,0.5);

  \draw[themegray, line width=0.7pt] (0.20,-0.012) -- (0.20,0.012);
  \draw[themegray, line width=0.7pt] (0.80,-0.012) -- (0.80,0.012);

  % 2. 函数主曲线（中层）
  \draw[themecurve, line width=2.0pt] (0,0) -- (1,1);
  \draw[themeamber, line width=2.0pt] (0,1) -- (1,0);

  % 3. 蛛网二周期迭代轨道（大箭头，高亮层）
  \draw[themealert, line width=1.8pt, -{Stealth[scale=1.4, length=3.2mm, width=2.4mm]}] (0.20,0) -- (0.20,0.80);
  \draw[themealert, line width=1.8pt, -{Stealth[scale=1.4, length=3.2mm, width=2.4mm]}] (0.20,0.80) -- (0.80,0.80);
  \draw[themealert, line width=1.8pt, -{Stealth[scale=1.4, length=3.2mm, width=2.4mm]}] (0.80,0.80) -- (0.80,0.20);
  \draw[themealert, line width=1.8pt, -{Stealth[scale=1.4, length=3.2mm, width=2.4mm]}] (0.80,0.20) -- (0.20,0.20);
  \draw[themealert, line width=1.8pt, -{Stealth[scale=1.4, length=3.2mm, width=2.4mm]}] (0.20,0.20) -- (0.20,0.80);

  % 4. 重点标记点
  \fill[themecurve] (0.5,0.5) circle (2.8pt);
  \fill[themealert] (0.20,0.80) circle (2.6pt);
  \fill[themealert] (0.80,0.20) circle (2.6pt);

  % 5. 所有文字注记与遮罩（顶层，绝不被任何线条穿过覆盖）
  \node[text=themeamber, fill=themebg, inner sep=0.8pt, anchor=south west] at (0.10,0.92) {$y=1-x$};
  \node[text=themecurve, fill=themebg, inner sep=0.8pt, anchor=south east] at (0.94,0.96) {$y=x$};
  \node[text=themecurve, fill=themebg, inner sep=0.8pt, anchor=north west] at (0.53,0.48) {$L=\dfrac12$};

  \node[text=themefg, fill=themebg, inner sep=0.8pt, anchor=north] at (0.20,-0.018) {$x_1=0.2$};
  \node[text=themefg, fill=themebg, inner sep=0.8pt, anchor=north] at (0.80,-0.018) {$x_2=0.8$};

  % 结论徽标
  \node[
    draw=themegray,
    rounded corners=1.5mm,
    fill=themebg,
    inner xsep=6pt,
    inner ysep=4pt,
    anchor=north,
    text=themefg
  ] at (0.50,-0.15) {唯一不动点仍不足以保证收敛：这里 $x_{n+2}=x_n$};
\end{scope}
\end{tikzpicture}
\end{CJK*}
\end{document}
